%%
% 结论
%%

\begin{conclusion}
\begingroup
\setcounter{section}{0}
\renewcommand{\thesection}{\arabic{section}}
\renewcommand{\theHsection}{conclusion.\arabic{section}}
\ctexset{section/number = \arabic{section}}

\section{主要工作与结论}

本文面向树莓派 5 这类无 GPU ARM 边缘平台上的 LeRobot 在线控制任务，
围绕 SO-101 机械臂、双摄像头输入和 ACT 策略模型，
研究端到端机器人策略在“观测--推理--执行”闭环中的系统级性能瓶颈。
与只统计端到端 FPS 或只从模型、外设单点优化入手的做法不同，
本文将问题拆解到 AI 推理框架层、Python 运行时层、Linux 内核层与硬件 I/O 层，
建立了面向机器人在线控制负载的跨层性能剖析方法。

在实验平台与测量方法方面，本文构建了可复现的 LeRobot 在线控制实验环境，
并结合 Python 阶段计时、strace 工具校准、ftrace 追踪和必要的内核插桩，
对主循环中的观测、推理、执行和等待阶段进行分层记录。
实验表明，strace 会显著放大上下文切换并改变阶段时间结构，
不适合承担定量时序分析；ftrace 延迟导出方案对主循环扰动相对可控，
更适合作为后续内核态观测工具。
在此基础上，本文进一步拆分了 pselect6、futex
等路径中的阻塞时间与 CPU 处理时间，使用户态日志和内核态事件能够互相印证。

系统级剖析结果表明，当前平台的主要瓶颈并不在摄像头读取、
舵机写入或单个 Linux 系统调用上，而在 ARM CPU 上执行 ACT 模型前向推理的绝对耗时。
摄像头路径已经通过 OpenCV 后台线程实现异步预取，主线程观测开销较低；
SCServo 舵机读写的毫秒级延迟主要受 1 Mbps 半双工总线和协议交互约束，
系统调用本身只占很小比例。相比之下，单次完整 ACT 推理达到秒级耗时，
其中 ResNet18 视觉编码和 Transformer Encoder 占据主要比例，
是限制在线控制频率的核心因素。

基于上述结论，本文进一步验证了异步推理对机器人在线控制的实际收益。
在 chunk\_size 为 100、目标 30 FPS 的实验设置下，
异步推理触发阈值取 0.30 时有效控制频率达到 27.42 FPS，
相对同步 baseline 提升 70.0\%，并在稳态阶段显著减少动作队列 stall。
这说明在不改变模型结构和硬件平台的条件下，
通过基础软件层面的流水线调度可以回收动作执行期间的空窗，
缓解长耗时推理对控制闭环的阻塞。

综上，本文证明了边缘机器人系统中的性能问题不能简单归因于某个模型、
某个外设或某次系统调用，而需要放在完整的软件栈和控制闭环中分析。
面向 LeRobot 这类机器人学习框架，系统级分层 profiling 不仅能够解释当前瓶颈的来源，
也能为后续推理运行时优化、操作系统调度改进和机器人基础软件架构设计提供依据。

\section{研究局限性}

本文的实验对象固定在树莓派 5、SO-101 机械臂和 ACT 策略模型上，
因此结论首先服务于 ARM CPU 边缘机器人这一类受限平台。
不同机械臂的通信协议、相机驱动、策略模型结构和任务动作尺度可能改变各阶段占比；
在模型更小、舵机总线更慢或相机分辨率更高的配置中，
瓶颈也可能从推理运行时转移到硬件 I/O 或图像预处理路径。
因此，本文给出的具体数值不应直接外推到所有机器人平台，
更适合作为一套跨层性能分析方法和基础软件优化思路的案例验证。

其次，本文已经对异步推理进行了实测验证，
但对推理运行时本体的优化仍然有限。
本文能够说明当前系统为什么慢、异步调度能恢复多少吞吐，
但还不能给出降低单次 ACT 推理耗时的最终工程方案。
此外，本文的动作质量评价仍以操作者现场观察为主，
尚未建立独立于主观判断的任务质量指标。
异步推理中的 stall 比例、过期帧比例和有效 FPS 可以反映控制流水线的运行状态，
但还不能完全替代抓取成功率、末端轨迹误差、动作平滑度和安全裕量等机器人任务指标。

\section{未来工作}

未来工作可以继续围绕机器人基础软件栈展开，并进一步上升到智能泛在操作系统的视角。
本文的实验说明，LeRobot 在线控制并不是单纯的模型推理问题，
而是由 AI 推理、反馈控制、硬件 I/O、用户态运行时和 Linux 调度共同构成的人机物融合软件系统。
后续可沿着“系统架构--任务调度--驱动生态”的主线，
将当前针对单一 LeRobot 平台的性能剖析扩展为面向边缘机器人设备的系统化构造方法。

在系统架构层面，可以围绕机器人实时控制负载探索更低开销的运行支撑机制。
当前系统仍建立在通用 Linux 宏内核和用户态 Python 运行时之上，
相机、串口、推理服务和控制循环需要通过多层软件抽象间接共享有限 CPU 资源。
后续可以以树莓派 5 这类 ARM 边缘平台为原型，
研究安全直通访问、统一内存视图和低开销用户态资源管理机制，
使相机缓冲区、图像张量、推理输入输出和执行器状态能够在受控权限下高效传递。

在任务与系统调度层面，需要从本文的单模型异步推理扩展到多模型协同推理和反馈控制调度。
后续可建立面向机器人在线控制的智能任务抽象，
描述视觉编码、动作策略、语言指令、安全监测和任务规划等模型之间的数据依赖、
时限约束、动作新鲜度要求和融合策略。
在无 GPU 平台上，重点仍应放在 ARM CPU 推理运行时与 Linux 调度策略的协同优化；
在具备 NPU、GPU 或多处理器资源的平台上，
则可进一步研究不同模型在异构计算资源之间的分配、迁移与决策融合。

在硬件 I/O 和驱动生态方面，后续仍需区分
“可以被软件隐藏的等待”和“由协议物理约束决定的等待”。
摄像头路径可继续从 V4L2 缓冲队列、OpenCV 后台线程、MJPEG 解码和图像张量转换等层面分析端到端时延；
舵机路径则应围绕 SCServo 半双工总线、串口波特率、同步读写协议和潜在的流水线读机制继续建模。
在此基础上，可进一步抽象不同机器人设备的操作原语，
形成面向传感器、相机、执行器和通信总线的统一接口与用户态驱动 SDK。

最后，本文使用的跨层 profiling 方法仍可进一步工具化，并与任务级行为质量评价结合。
后续可以将 Python 日志、torch.profiler、ftrace、内核自定义插桩和实验 CSV
自动对齐为统一时间线，形成面向机器人基础软件的可复用性能诊断工具。
同时，还需要将系统级性能指标与抓取成功率、末端轨迹误差、动作平滑度、
安全裕量和故障降级次数等任务级指标关联起来，
避免只优化 FPS 而忽略机器人行为质量。
\endgroup
\end{conclusion}
